% formal/stateless.tex
% mainfile: ../perfbook.tex
% SPDX-License-Identifier: CC-BY-SA-3.0

\section{无状态模型检查器}
\label{sec:formal:Stateless Model Checkers}
%
\epigraph{他列了一个清单，对它做了两次排列\dots}
	{改编自 Haven Gillespie 和 J. Fred Coots}

上一节描述的 SAT solver 方法非常方便且强大，但对所有可能执行
（包括状态）的完整跟踪可能会产生大量开销。
事实上，内存和 CPU 时间开销可能会严重限制可行验证的程序规模，
这引发了一个问题：不那么精确的方法是否可能在更大的程序中找到 bug。

\begin{figure}
\centering
\resizebox{2.1in}{!}{\includegraphics{formal/nidhugg}}
\caption{Nidhugg 处理流程}
\label{fig:formal:Nidhugg Processing Flow}
\end{figure}

虽然对这个问题的评判尚无定论，但 stateless model checker
（如 \IX{Nidhugg}~\cite{CarlLeonardsson2014Nidhugg}）在某些
情况下已经能处理更大的程序~\cite{SMC-TreeRCU}，并且具有
类似的易用性，如\cref{fig:formal:Nidhugg Processing Flow}所示。
此外，在某些 Linux 内核 RCU 验证场景中，Nidhugg 比
\co{cbmc} 快了一个数量级以上。
当然，Nidhugg 的速度和可扩展性优势与它不处理数据非确定性有关，
但这在这些特定的验证场景中不是一个因素。

尽管如此，与 \co{cbmc} 一样，Nidhugg 尚未能找到 Linux 内核 RCU
维护者之前不知道的 bug。
然而，它确实能够证明 Linux 内核 RCU 中的一个历史 bug
是由不同于维护者所认为的提交修复的，
这给了人们更多的希望，即像 Nidhugg 这样的 stateless model checker
有朝一日可能对在并行代码中找到并发 bug 有所帮助。
